% Part: propositional-logic
% Chapter: syntax-and-semantics
% Section: introduction

\documentclass[../../../include/open-logic-section]{subfiles}

\begin{document}

\olfileid{pl}{syn}{int}

\olsection{Inleiding}

In de propositielogica gaat het om !!{formula}s. We bouwen die op uit
!!{propositional variable}s met de verbindingstekens $\lnot$, $\land$,
$\lor$, $\lif$ en $\liff$. Je kunt !!a{propositional variable}~$p$
zien als een letter die staat voor een zin of bewering die waar of
onwaar is. Als de ``waarheidswaarde'' van de !!{propositional
variable}s in !!a{formula} bekend is, ligt ook de waarheidswaarde vast
van elke !!{formula} die we daaruit met de verbindingstekens opbouwen.
Daarom heet de propositielogica \emph{waarheidsfunctioneel}: haar
betekenisregels werken met functies van waarheidswaarden. Andere
aspecten van waarheid en onwaarheid laten we buiten beschouwing. Het
maakt hier bijvoorbeeld niet uit of iets noodzakelijk waar is en niet
alleen toevallig waar, of iemand weet dat het waar is, of het nu waar
is in plaats van vroeger of later. We gebruiken maar twee waarden:
waar ($\True$) en onwaar ($\False$). We bespreken dus niet de
mogelijkheid dat een bewering noch waar noch onwaar is, of maar half
waar. Ook nemen we alleen verbindingstekens waarbij de
waarheidswaarden van de delen de waarheidswaarde van de hele
!!{formula} volledig bepalen, dus niet bijvoorbeeld de betekenis van
die delen. Bij voorwaardelijke zinnen in het Engels is omstreden
of hun waarheidswaarde op deze manier wordt bepaald. Voor de materiële
implicatie~$\lif$ gebeurt dat wel. Andere logica's behandelen
voorwaardelijke verbindingen waarbij dat niet zo is.

Voordat we de theorie zelf en de metatheorie, waarin we die theorie
onderzoeken, kunnen uitwerken, moeten we twee dingen vastleggen: hoe de
uitdrukkingen worden gevormd en wat ze betekenen. Dat zijn de
vormregels (syntaxis) en de betekenisregels (semantiek). We geven één
manier om met verbindingstekens !!{formula}s uit !!{propositional
variable}s op te bouwen. Andere definities zijn ook mogelijk. Een
ander systeem kan andere symbolen kiezen, andere verbindingstekens als
basis nemen en haakjes anders gebruiken. Het kan de haakjes zelfs
helemaal weglaten, zoals in de zogenoemde Poolse notatie. In al die
benaderingen definiëren de vormingsregels de verzameling !!{formula}s
\emph{inductief}. Als de regels goed zijn opgesteld, kan elke
uitdrukking volgens die regels in wezen maar op één manier zijn
ontstaan. Elke uitdrukking is dan \emph{uniek leesbaar}. Daardoor
kunnen we de betekenis op dezelfde stapsgewijze, inductieve manier
vastleggen.

De betekenis van uitdrukkingen hoort bij de semantiek. Daar staat de
vraag centraal of een formule waar is onder !!a{valuation}. Zo'n
!!{valuation}~$\pAssign{v}$ kent aan elke !!{propositional variable}
  de waarde $\True$ of $\False$ toe. Daarmee bepaalt de !!{valuation}
  de waarde $\pValue{v}(!A)$ voor elke !!{formula}~$!A$. !!^a{formula}
wordt door !!a{valuation}~$\pAssign{v}$ vervuld precies als
$\pValue{v}(!A) = \True$. Dat schrijven we als $\pSat{v}{!A}$. Ook
deze relatie kunnen we stap voor stap definiëren aan de hand van de
opbouw van~$!A$. Met de waarheidsfuncties van de logische
verbindingstekens bepalen we bijvoorbeeld wanneer $!A \land !B$ wordt
vervuld: dat hangt ervan af of $!A$ en~$!B$ wel of niet worden
vervuld.

Met de relatie $\pSat{v}{!A}$ voor zinnen kunnen we daarna drie
belangrijke begrippen definiëren: tautologie, semantisch gevolg en
vervulbaarheid. !!^a{formula} is een tautologie, $\Entails !A$, als
  elke !!{valuation} haar vervult. Anders gezegd:
  $\pValue{v}(!A) = \True$ voor elke~$\pAssign{v}$. Een formule volgt uit een
verzameling !!{formula}s, $\Gamma \Entails !A$, als elke
!!{valuation} die alle !!{formula}s in~$\Gamma$ vervult, ook~$!A$
vervult. Een verzameling !!{formula}s is vervulbaar als er een
!!{valuation} is die alle formules in die verzameling tegelijk
vervult. Zowel de !!{formula}s als hun vervulling zijn inductief
gedefinieerd. Daarom kunnen we met inductie eigenschappen van deze
betekenisregels bewijzen en de begrippen met elkaar verbinden.
\end{document}
